% Review draft: FF.46. Generic-site converse and all-prime Kummer rigidity.
\subsection{Cyclic rational covers determine every inertia exponent}

\begin{claim}{FF.46}{Proved}\label{clm:FFallprimekernel}
Let \(k\) be a characteristic-zero field and \(X/k\)
a smooth geometrically integral curve with function field
\(F\). The following data determine an automorphism of
\(X_{\mathrm{\acute et}}\) over \(BG_k\): a continuous
automorphism \(\alpha\) of \(G_F\) over \(G_k\),
and a \(G_k\)-equivariant bijection
\(\theta:X(\bar k)\to X(\bar k)\), such that
\(\alpha\) sends the geometric conjugacy class of
\(I_x\) to that of \(I_{\theta(x)}\) for every
\(x\in X(\bar k)\). The resulting actual point
function is \(\theta\), and the resulting generic
action is \(\alpha\), with the usual geometric inner
ambiguity.

Now let \(E/B_p\) be finite, and let \(\varphi\)
be a pointwise-trivial pinned automorphism of
\(\mathbb A^1_{E,\mathrm{\acute et}}\) over \(BG_E\).
For every finite extension \(L/E\) and every positive
integer \(m\), its generic action fixes
\[
                    H^1(G_{L(T)},\mu_m)
\]
pointwise. Equivalently, it acts trivially on the maximal
abelian profinite quotient of \(G_{L(T)}\), after
every finite constant base change. All geometric inertia
exponents are one, at every prime. Neither assertion
identifies the remaining nonabelian generic action with
an inner automorphism.
\end{claim}

\begin{proof}
\emph{Reconstructing the generating site.}
We first give a converse to the construction in
Claim~\ref{clm:FFgenerickernel}. Fix prolongations and
geometric residue embeddings at the closed points of
\(X\). The category of quasicompact separated etale
\(X\)-schemes has the following description:
\[
 \begin{gathered}
 \left(S,\ (B_x)_{x\in|X|_{\rm cl}}\right),
 \qquad S\text{ a finite continuous }G_F\text{-set},\\
 B_x\subset S^{I_x}\text{ a }D_x\text{-stable subset}
       \qquad(x\in|X|_{\rm cl}).
 \end{gathered}
                                                        \tag{FF.46.1}
\]
where \(B_x=S^{I_x}\) for all but finitely many closed
\(x\). Its morphisms are the \(G_F\)-maps on generic
sets that take every \(B_x\) into its counterpart.
Indeed, form the finite normalization belonging to
\(S\), take its etale locus \(j_*S\), and remove
the finitely many closed points specified by the omitted
\(D_x\)-orbits. This reconstructs the object. Conversely,
Claim~\ref{clm:FFgenerickernel} identifies every object
as precisely such an open. A generic map extends to the
finite normalizations, and its image lies in the indicated
open exactly when the conditions on the \(B_x\) hold.
This proves the asserted description of arrows as well.
Multiple generic components are allowed in \(S\).
Nonseparated objects are not needed: this separated
subcategory is the generating site by
\citet[Theorem~A.8 and Corollary~A.11]{FF_CHW2024}.

Let \(N=G_{\bar k(X)}\). For a chosen geometric
\(x\), choose \(g_x\in N\) such that
\[
                  \alpha(I_x)=g_x I_{\theta(x)}g_x^{-1}.
                                                        \tag{FF.46.2}
\]
On generic sets use pullback by \(\alpha\), so the new
action on \(S\) is \(g\cdot_{\rm new}s=\alpha(g)s\).
Replace the local subset at \(x\) by
\[
                                g_x B_{\theta(x)}.
                                                        \tag{FF.46.3}
\]
This is a subset of \((\alpha^*S)^{I_x}\). It is
\(D_x\)-stable because (FF.45.4) and (FF.46.2) imply
\(\alpha(D_x)=g_xD_{\theta(x)}g_x^{-1}\).
Another choice of \(g_x\) differs on the right by an
element of \(D_{\theta(x)}\), again by (FF.45.4), and
therefore gives the same subset. Changes of valuation
prolongations give the usual identical description up
to canonical generic-set isomorphism. The
\(G_k\)-equivariance of \(\theta\) makes the rule
well-defined on closed-point orbits. Only finitely many
full subsets are changed, since \(\theta\) bijects
those orbits.

The rule preserves arrows and has an inverse obtained
from \(\alpha^{-1},\theta^{-1}\). Finite limits are
computed on the generic and closed geometric stalks, as
are jointly surjective etale covering families. Thus it
is an equivalence of the generating sites and induces
an equivalence of their topoi. On sets inflated from
\(G_k\), the generic action is the specified identity
and \(g_x\in N\) acts trivially. Hence the equivalence
has the required structure over \(BG_k\).

For completeness, the residue-field identification is
also determined. Since \(\theta\) is an equivariant
bijection, the stabilizers of \(x\) and \(\theta(x)\)
in \(G_k\) coincide. Their embedded finite residue
fields inside \(\bar k\) are consequently the same.
The isomorphism on \(D_x/I_x\) has the identity
projection to this stabilizer, since both \(\alpha\)
and the conjugation in (FF.46.2) are relative over
\(G_k\). The stalk identifications in (FF.46.3)
therefore identify the geometric point functor at \(x\)
with the target point functor at \(\theta(x)\), with
its specified base comparison. By
\citet[Proposition~3.6]{FF_CHW2024}, the actual point
function is exactly \(\theta\). The generic action
is the prescribed one by construction. This proves the
first assertion.

\emph{The cyclic rational cover.}
Fix now \(L/E\) finite and \(a\in L\), and base
change \(\varphi\) to \(L\). Its point function is
still the identity. On
\(V=\mathbb A^1_L\setminus\{a\}\), geometric first
cohomology with \(\mu_m\)-coefficients is generated
by \(\kappa_m(T-a)\). The equivalence therefore gives
\[
 \varphi^*\kappa_m(T-a)
       =u_m\kappa_m(T-a)+\eta_m,
 \quad u_m\in(\mathbb Z/m\mathbb Z)^\times,
 \quad\eta_m\in H^1(G_L,\mu_m).
                                                        \tag{FF.46.4}
\]
The displayed geometric generator is already lifted by
the indicated global unit, so the remaining term is
constant, by inflation--restriction. Evaluation at
\(T=a+1\) gives \(\eta_m=0\).

The field \(L\) contains \(\mu_m\). Thus the
extension
\[
              L(T)\subset L(U),\qquad T=U^m+a,
                                                        \tag{FF.46.5}
\]
is cyclic Galois of degree \(m\). Its absolute Galois
subgroup is the kernel of the homomorphism
\(\kappa_m(T-a):G_{L(T)}\to\mu_m\).
By (FF.46.4), the generic automorphism \(\alpha\)
preserves this subgroup. Its restriction
\(\alpha_U\) is an automorphism of \(G_{L(U)}\)
over \(G_L\).

The exact local theory in
\citet[Theorem~4.1 and Corollaries~4.3--4.5]
{FF_Koenigsmann2003} applies to \(L\), and after every
finite constant extension, both for \(\mathbb P^1_T\)
and \(\mathbb P^1_U\). It gives the associated
\(G_L\)-equivariant permutation \(\theta_U\) of
geometric places. More explicitly, first make a chosen
geometric place rational over a finite constant field
\(L'/L\). Corollary~4.3 characterizes its arithmetic
decomposition conjugacy class. Any conjugating element
can be multiplied by the inverse of an element of the
target decomposition group with the same projection to
\(G_{L'}\), because that projection is surjective.
The resulting conjugator is geometric. Applying this
argument over all finite constants, uniqueness of the
recovered valuation classes makes the point permutations
compatible with scalar extension and with the
\(G_L\)-action. This is the geometric conjugacy
transport required in (FF.46.2).
This permutation lies above the
identity permutation on \(\mathbb P^1_T\). To see
the needed compatibility explicitly, choose a place of
the common separable closure. Its decomposition and
inertia groups for \(L(U)\) are the intersections of
the corresponding groups for \(L(T)\) with
\(G_{L(U)}\). The subgroup is preserved by
\(\alpha\), so these intersections and their place
restrictions are transported compatibly. Uniqueness in
the local characterization identifies the transported
places with those of \(\theta_U\). In particular,
\[
                           \theta_U(u)^m=u^m
                  \qquad(u\in\bar L).
                                                        \tag{FF.46.6}
\]
The unique place over \(T=\infty\) is fixed, as is
the unique place \(U=0\) over \(T=a\). Therefore
\(\alpha_U,\theta_U\) satisfy the first assertion
for the affine line \(\mathbb A^1_U\). They define
a pinned automorphism of its full etale topos over
\(BG_L\), rather than merely of its generic finite
cover category.

Claim~\ref{clm:FFmetabelianaffine} now applies to this
lifted automorphism. Its actual point function is affine.
Since it fixes zero and satisfies (FF.46.6), it has the
form
\[
                       \theta_U(u)=c u,
                       \qquad c\in\mu_m(L).
                                                        \tag{FF.46.7}
\]
On the other hand, (FF.46.4) says that \(\alpha\)
acts on the cyclic deck quotient of (FF.46.5) by the
power \(u_m\). Conjugation and transport of inertia
classes are compatible. Indexing the deck rotations by
their action on the geometric points of the cover,
this gives
\[
                    \theta_U(\zeta u)
                          =\zeta^{u_m}\theta_U(u)
                    \qquad(\zeta\in\mu_m).
                                                        \tag{FF.46.8}
\]
Combining
(FF.46.7) and (FF.46.8) gives
\(\zeta^{u_m}=\zeta\) for every \(\zeta\in\mu_m\).
Hence \(u_m=1\) in \(\mathbb Z/m\mathbb Z\).
We have proved, for every \(a\) and \(m\),
\[
                  \varphi^*\kappa_m(T-a)=\kappa_m(T-a).
                                                        \tag{FF.46.9}
\]

\emph{Descent of an arbitrary rational Kummer class.}
For \(r\in L(T)^\times\), take its nonvanishing
open \(V_r\subset\mathbb A^1_L\) and a finite
Galois constant extension \(L'/L\) splitting its
divisor. Over \(L'\), it is a constant times a product
of integer powers of the linear units in (FF.46.9).
The relative structural isomorphism fixes the constant
classes. Thus
\(\delta=\varphi^*\kappa_m(r)-\kappa_m(r)\)
restricts to zero on \((V_r)_{L'}\).
Inflation--restriction for this purely constant cover
puts \(\delta\) in the image of
\(H^1(\operatorname{Gal}(L'/L),\mu_m)\), hence in
the constant subgroup from \(H^1(G_L,\mu_m)\).
Choose \(t_0\in V_r(L)\). Actual pointwise identity
gives \(t_0^*\delta=0\), and evaluation at \(t_0\)
is a left inverse on constant classes. Therefore
\(\delta=0\) before the scalar extension as well.

Every element of \(H^1(G_{L(T)},\mu_m)\) is a
Kummer class. Moreover \(\mu_m\subset L\), so
these cohomology groups are exactly the continuous
characters with values in \(\mu_m\). Their equality
for every \(m\) is equivalent to equality on the full
abelian profinite quotient. The exponents on finite
geometric inertia are detected by (FF.46.9), and the
single peripheral relation on the projective line gives
the same exponent at infinity. All of them are one.
\end{proof}

This improves the direct \(p\)-primary calculation:
the other primes are controlled by the rational cyclic
covers and the already-proved affine point-function
rigidity, even though point Kummer classes at those
primes vanish over \(B_p\). A nonabelian generic
automorphism acting trivially on all these abelian
quotients has not been excluded.
